\section{The actual relative-to-ordinary image on the full finite packet spectrum}
\label{sec:packet-interior}

The finite packet spectrum in TP8 has several arithmetic closed
points. Its ordinary cohomology can therefore be calculated directly,
rather than inferred from the irreducible arithmetic base. Let
$P=\SpecLat L$ be the full finite nonempty support poset. Let $Z_\R$
be the finite discrete spectrum of the nonempty packet algebra
$A_{Z,\R}$ of TP1, and put $c=|Z_\R|\ge1$. A critical root counts once and an
off-critical reflected pair counts once, by TP5. Its full split
spectrum is exactly
\[
 Y_Z=P\triangleright Z_\R.
 \tag{PI1}\label{pi:space}
\]
We use the constant sheaf with value a specified complex vector
space $W$ and all its restriction maps equal to the identity.
This is a sheaf of complex vector spaces on the actual prime
topology. It is not silently made a module over the semiring
structure sheaf. At every point its stalk is the stated $W$;
this choice is part of the construction.

Write $C^\bullet(P;W)$ for the strict-chain cochain complex,
with its alternating face differential $\delta$, and augment it
by $W$ in degree $-1$ with the constant-section map. Denote this
augmented complex by $\widetilde C^\bullet(P;W)$. The full
finite-space resolution used in FL14--19 computes constant-sheaf
cohomology by these cochains. The relative complex for the closed
arithmetic part is the kernel of restriction from $Y_Z$ to $P$.
Every strict chain of $Y_Z$ either lies entirely in $P$ or ends
at exactly one arithmetic point $z$. Thus, for $n\ge1$,
\[
 C^n(Y_Z;W)=C^n(P;W)\oplus
       \bigoplus_{z\in Z_\R}C^{n-1}(P;W),
 \quad
 C^0(Y_Z;W)=C^0(P;W)\oplus W^{Z_\R}.
 \tag{PI2}\label{pi:chains}
\]
For $(a,(b_z)_z)$ in degree $n$, its differential has components
\[
 \delta a,\qquad
 \delta b_z+(-1)^{n+1}a.
 \tag{PI3}\label{pi:differential}
\]
For $n=0$, $\delta b_z$ means the constant cochain with value
$b_z$, so the second component is $b_z-a(x)$ on the edge $(x,z)$.
For all other degrees the formula follows by deleting each
vertex of $(x_0,\ldots,x_n,z)$ with its actual alternating sign.
This derives the differential, including the connecting sign.

\begin{theorem}[Complete relative image and its kernel]
For every $n\ge1$ there are canonical identifications
\[
 \begin{aligned}
 H^n_{Z_\R}(Y_Z;W)&=\C^{Z_\R}\otimes_\C
                         \widetilde H^{n-1}(P;W),\\
 H^n(Y_Z;W)&=(\C^{Z_\R}/\Delta\C)\otimes_\C
                         \widetilde H^{n-1}(P;W).
 \end{aligned}
 \tag{PI4}\label{pi:groups}
\]
The relative-to-ordinary map is the displayed quotient on the
first factor. It is onto, and its kernel is
$\Delta\C\otimes\widetilde H^{n-1}(P;W)$.
Also $H^0_{Z_\R}(Y_Z;W)=0$ and $H^0(Y_Z;W)=W$.
\end{theorem}
\begin{proof}
By PI2--3, the kernel of restriction has in degree $n$ the
direct sum of $c$ copies of $\widetilde C^{n-1}(P;W)$, with
differential $\delta$. It computes the first group of PI4.
The constant augmentation is injective because $P\ne\varnothing$,
so its degree-zero relative group vanishes. The connecting map
from $H^n(P;W)$ sends $[a]$ to
$((-1)^{n+1}[a])_{z\in Z_\R}$ in the relative degree $n+1$.
For $n=0$ this is the negative diagonal map after quotienting constant
sections, and for $n>0$ it is a signed diagonal injection.
The long exact sequence of the surjective cochain restriction
therefore gives the second group of PI4 and the exact quotient
map, with no choice of an arithmetic point. In degree zero,
every edge forces all values to coincide; the poset is connected
because $P$ and $Z_\R$ are nonempty and every pair $(x,z)$ is an
edge. Its global constant sections are precisely $W$.
\end{proof}

The ordinary quotient also has the following explicit cochain
receiver. Put $U=\C^{Z_\R}$, $V=U/\Delta\C$ and
$A^\bullet=\widetilde C^\bullet(P;W)$. The map
\[
 \rho_n(a,b)=(q\otimes1)b:\ C^n(Y_Z;W)\longrightarrow
                      V\otimes A^{n-1},\qquad q:U\twoheadrightarrow V
 \tag{PI7}\label{pi:cochain-quotient}
\]
commutes with the differential because $q\Delta=0$. Its kernel
is $K^n=C^n(P;W)\oplus A^{n-1}$, with
$d_K(a,b)=(\delta a,\delta b+(-1)^{n+1}a)$.
For $n\ge1$, let $h_n(a,b)=((-1)^n b,0)$.
Substitution gives $d_Kh+hd_K=1$ in positive degrees.
In degree zero, with $\varepsilon(a,b)=b$ and
$\iota(w)=((w)_{x\in P},w)$, it gives
$h_1d_K=1-\iota\varepsilon$. Thus the kernel retracts to
$W[0]$. This proves the stated quotient on cohomology directly.
Its restriction to the relative complex is the same $q\otimes1$,
without a choice of an arithmetic point.

For Boolean support with $d$ atoms, $P$ consists of $d$ points.
The calculation specializes to
\[
 \dim_\C H^1_{Z_\R}(Y_Z;\C)=c(d-1),\qquad
 \dim_\C H^1(Y_Z;\C)=(c-1)(d-1),
 \qquad H^{n\ge2}(Y_Z;\C)=0.
 \tag{PI5}\label{pi:boolean}
\]
For the six-point triangle-boundary support, $\widetilde H^1(P;\C)
=\C$ and all other reduced groups vanish. Thus its ordinary
degree-two group is $\C^{Z_\R}/\Delta\C$, and its relative
degree-two group is $\C^{Z_\R}$. In particular this packet
has a nonzero relative-to-ordinary image when $c>1$. The example
comes from the same support complex in BP12, not an assigned
cohomological degree.

\subsection{Every original prime acts on these exact groups}
Take $W=A_Z$, the entire original complex jet algebra, and let
$\Pi_p$ act by multiplication by its full $p^s$ jet at every
stalk. This is a sheaf endomorphism: all restriction maps are
identities and commute with it. Formula PI3 proves chain
compatibility and gives the action $I\otimes\Pi_p$ under PI4.
If $h_{n-1}=\dim_\C\widetilde H^{n-1}(P;\C)$, then
\[
 \begin{aligned}
 \operatorname{Tr}\bigl(\Pi_p\mid H^n_{Z_\R}(Y_Z;A_Z)\bigr)
 &=c\,h_{n-1}\sum_{\rho\in Z}m_\rho p^\rho,\\
 \operatorname{Tr}\bigl(\Pi_p\mid H^n(Y_Z;A_Z)\bigr)
 &=(c-1)h_{n-1}\sum_{\rho\in Z}m_\rho p^\rho.
 \end{aligned}
 \tag{PI6}\label{pi:traces}
\]
Indeed the coefficient matrix on each full local jet factor is
lower triangular in the original ascending basis
$1,t_\rho,\ldots,t_\rho^{m_\rho-1}$, with diagonal $p^\rho$
repeated $m_\rho$ times. Explicitly its row $k$, column $j$ entry is
$p^\rho(\log p)^{k-j}/(k-j)!$ for $k\ge j$ and zero otherwise.
Tensoring with the explicitly counted topological factors gives
the displayed traces. Its nilpotent coefficients are still the
original Taylor coefficients; only the trace omits their strictly
triangular entries. The ordinary dual of
$\C^{Z_\R}/\Delta\C$ is the sum-zero subspace of
$(\C^{Z_\R})^*$, so the dual receiving map is also specified
without choosing a branch. Write $H=\widetilde H^{n-1}(P;\C)$,
and let $T$ be either $U\otimes H$ or $V\otimes H$. The exact
perfect dual receiver, using FJ4--5, is
\[
 1_{T^*}\otimes D_R:\ T^*\otimes\overline{A_Z}
     \xrightarrow{\sim}(T\otimes A_Z)^*,\qquad
 (\phi\otimes\bar f,t\otimes h)\longmapsto\phi(t)R(f,h).
 \tag{PI8}\label{pi:dual-receiver}
\]
On the target the prime acts by
$p(1\otimes\Pi_p^{-1})^*$; on the source it acts by
$1\otimes\overline{\Pi_p}$. The equality follows by applying
FJ6 to $R(\Pi_pf,h)=pR(f,\Pi_p^{-1}h)$, so every factor and
inverse is specified. The dual relative map is
$q^*\otimes1_{H^*}\otimes1_{A_Z^*}$, where $q^*$ embeds the
sum-zero covectors into $U^*$. No topological self-pairing is
assumed. Contracting by the original $J_g$ gives exactly
\[
 \phi(t)\sum_{\rho\in Z}m_\rho
       \overline{f(1-\bar\rho)}h(\rho).
 \tag{PI9}\label{pi:contracted-duality}
\]
The two radicals are $T^*\otimes\overline{\ker J_g}$ and
$T\otimes\ker J_g$, because the topological evaluation pairing
is perfect and FJ4--5 identifies the coefficient radical.
This retains all multiplicities through the exact dual maps.

Finally PI1 is the finite packet base, with the dense embedding
TP8 into $P\triangleright\Spec\mathcal B$. On the latter
irreducible arithmetic base, the constant-coefficient sheaf has
ordinary cohomology $\C[0]$: the sheaf-triple computation FL1--13
reduces it to the constant sheaf on $\Spec\mathcal B$, whose
restriction maps between nonempty opens are identities and
which is flasque. The ring $\mathcal B$ is an integral domain,
since after scalar extension it is the Laurent polynomial
domain $\C[X_2^{\pm1},X_3^{\pm1}]$. The same argument with
$A_Z$ coefficients gives $A_Z[0]$. Consequently the pullback
from that ordinary positive-degree cohomology into PI4 is the
zero map. The nonzero packet image calculated here has the exact
base-change map just stated; it has not been identified with
an already nonzero global interior group.

\begin{figure}[ht]
\centering
\begin{tikzcd}[column sep=large,row sep=large]
 \C^{Z_\R}\otimes\widetilde H^{n-1}(P;W)
 \arrow[r,two heads]
 &(\C^{Z_\R}/\Delta\C)\otimes\widetilde H^{n-1}(P;W)\\
 H^n_{Z_\R}(Y_Z;W)\arrow[u,equal]\arrow[r]
 &H^n(Y_Z;W)\arrow[u,equal]
\end{tikzcd}
\caption{The actual relative-to-ordinary map on the full finite
packet spectrum, PI2--4. Every support chain and every real
arithmetic closed point occurs. With the original jet coefficient
space, the prime action and trace on both sides are PI6.
No extra prime-weight factor is inferred from the degree shift.}
\end{figure}
